% GENERATED FILE. Edit canonical lessons, then run npm run generate:books.
% canonical-source-hash: fnv1a64:7869caca45a86b26
% canonical-lessons: BN-C275-nirbhor-kore, BN-C275-jetam, BN-C275-songge, BN-W275-letter-uma, BN-W275-six-lines, BN-R275-depends-and-would

\chapter{It Depends, and What Would Be}
\label{ch:bn-it-depends-and-what-would-be}
\hlchaptermodality{275}

\begin{chapteropening}
\textbf{By the end of this chapter:} \emph{I can say what something depends on (… -er opôr nirbhôr kôre), say what would happen and what used to happen (jetām), write six connected lines about a trip, and read and write \bn{ঙ} in \bn{সঙ্গে}.}

Complete BN-R275-depends-and-would: say what something depends on and what would be, after writing six connected lines about a trip.
\end{chapteropening}

\section[\texorpdfstring{… -\bn{এর} \bn{ওপর} \bn{নির্ভর} … (… -er opôr …)}{… -এর ওপর নির্ভর … (… -er opôr …)}]{… -\bn{এর} \bn{ওপর} \bn{নির্ভর} \bn{করে} (-er opôr nirbhôr kôre) — it depends on}

\label{lesson:BN-C275-nirbhor-kore}



\begin{quote}
Say \textquotedblleft{}If it rains, we'll stay at home\textquotedblright{}. (\emph{brishṭi hôle āmrā bāṛite thākbô}.) Then \textbf{\bn{ওপরে}}, \textquotedblleft{}above, upstairs\textquotedblright{}.
\end{quote}

\subsection*{You'll want to know: it depends on}
\begin{quote}
\textbf{\bn{এটা} \bn{আবহাওয়ার} \bn{ওপর} \bn{নির্ভর} \bn{করে।}}
\end{quote}

\begin{quote}
\emph{eṭā ābôhāoyār opôr nirbhôr kôre.} — \textquotedblleft{}It depends on the weather.\textquotedblright{}
\end{quote}

\textbf{\bn{ওপর}} (\emph{opôr}), \textquotedblleft{}on, upon\textquotedblright{}, is the word inside \textbf{\bn{ওপরে}}. \textbf{\bn{নির্ভর} \bn{করা}} (\emph{nirbhôr kôrā}) is \textquotedblleft{}to depend\textquotedblright{}. What it depends on takes the \textbf{-\bn{র}} form, then \textbf{\bn{ওপর}}: \textquotedblleft{}it does depending upon the weather\textquotedblright{}. \textbf{\bn{সব}}, \textquotedblleft{}all\textquotedblright{}, from the last chapter, makes it \textquotedblleft{}everything\textquotedblright{}: \textbf{\bn{সব} \bn{আবহাওয়ার} \bn{ওপর} \bn{নির্ভর} \bn{করে}}.

\subsection*{You'll want to know: it's up to you}
With a person, it is \textquotedblleft{}up to\textquotedblright{} them: \textbf{\bn{এটা} \bn{তোমার} \bn{ওপর} \bn{নির্ভর} \bn{করে}} (\emph{eṭā tomār opôr nirbhôr kôre}), \textquotedblleft{}it's up to you\textquotedblright{}. About a plan:

\begin{quote}
\textbf{\bn{কাল} \bn{বাজারে} \bn{যাব}, \bn{যদি} \bn{মীরা} \bn{আসে।} \bn{সব} \bn{মীরার} \bn{ওপর} \bn{নির্ভর} \bn{করে।}}
\end{quote}

\begin{quote}
\emph{kāl bājāre jābô, jôdi mīrā āse. sôb mīrār opôr nirbhôr kôre.} — \textquotedblleft{}We'll go to the market tomorrow if Mira comes. Everything depends on Mira.\textquotedblright{}
\end{quote}

\subsection*{Your turn}
Say these aloud:

\begin{itemize}
\raggedright
  \item \textquotedblleft{}It depends on the weather\textquotedblright{} — \emph{eṭā ābôhāoyār opôr nirbhôr kôre}
  \item \textquotedblleft{}It's up to you\textquotedblright{} — \emph{eṭā tomār opôr nirbhôr kôre}
  \item \textquotedblleft{}Everything depends on the train\textquotedblright{} — \emph{sôb ṭrener opôr nirbhôr kôre}
\end{itemize}

\begin{tcolorbox}[breakable,colback=teal!4,colframe=teal!35!black,title={Before you move on}]
\textquotedblleft{}It depends on the weather\textquotedblright{}? (\textbf{\bn{এটা} \bn{আবহাওয়ার} \bn{ওপর} \bn{নির্ভর} \bn{করে}}.) \textquotedblleft{}It's up to you\textquotedblright{}? (\textbf{\bn{তোমার} \bn{ওপর} \bn{নির্ভর} \bn{করে}}.)
\end{tcolorbox}
\section[\texorpdfstring{\bn{যদি} … \bn{থাকত} … (jôdi … thāktô …)}{যদি … থাকত … (jôdi … thāktô …)}]{\bn{যদি} … \bn{থাকত}, \bn{তাহলে} … \bn{যেতাম} (jôdi … thāktô, tāhôle … jetām) — what would be}

\label{lesson:BN-C275-jetam}



\begin{quote}
Say \textquotedblleft{}It depends on the weather\textquotedblright{}. (\emph{eṭā ābôhāoyār opôr nirbhôr kôre}.) Then \textquotedblleft{}If you come, I'll cook\textquotedblright{}. (\emph{jôdi tumi āso, āmi rānnā kôrbô}.) That if can come true.
\end{quote}

\subsection*{You'll want to know: what would be}
\begin{quote}
\textbf{\bn{যদি} \bn{সময়} \bn{থাকত}, \bn{তাহলে} \bn{আমি} \bn{যেতাম।}}
\end{quote}

\begin{quote}
\emph{jôdi sômôy thāktô, tāhôle āmi jetām.} — \textquotedblleft{}If I had time, I would go.\textquotedblright{}
\end{quote}

For an if that is not true, both halves take the \textbf{-\bn{ত}-} form: the stem, then \textbf{-\bn{ত}-}, then the person: \textbf{\bn{আমি} \bn{যেতাম}} (\emph{āmi jetām}), \textbf{\bn{তুমি} \bn{যেতে}} (\emph{tumi jete}), \textbf{\bn{সে} \bn{যেত}} (\emph{she jetô}), \textbf{\bn{আপনি} \bn{যেতেন}} (\emph{āpni jeten}). \textbf{\bn{যাওয়া}} $\to$ \textbf{\bn{যে}-}, \textbf{\bn{খাওয়া}} $\to$ \textbf{\bn{খে}-}, \textbf{\bn{হওয়া}} $\to$ \textbf{\bn{হ}-}; \textbf{\bn{করা}} $\to$ \textbf{\bn{করতাম}}, \textbf{\bn{থাকা}} $\to$ \textbf{\bn{থাকতাম}}. \textbf{\bn{টাকা} \bn{থাকলে} \bn{আমি} \bn{একটা} \bn{নৌকা} \bn{কিনতাম}} (\emph{ṭākā thākle āmi ekṭā noukā kintām}), \textquotedblleft{}if I had money, I'd buy a boat\textquotedblright{}.

\subsection*{You'll want to know: what used to be}
The same form tells what happened again and again in the past:

\begin{quote}
\textbf{\bn{আগে} \bn{আমি} \bn{রোজ} \bn{নদীতে} \bn{যেতাম।}}
\end{quote}

\begin{quote}
\emph{āge āmi roj nôdīte jetām.} — \textquotedblleft{}I used to go to the river every day.\textquotedblright{}
\end{quote}

\textbf{\bn{আগে}}, \textquotedblleft{}before\textquotedblright{}, here is \textquotedblleft{}in the old days\textquotedblright{}. A word like \textbf{\bn{রোজ}} or \textbf{\bn{প্রতিদিন}} makes the habit clear.

\subsection*{Your turn}
Say these aloud:

\begin{itemize}
\raggedright
  \item \textquotedblleft{}If I had time, I would go\textquotedblright{} — \emph{jôdi sômôy thāktô, tāhôle āmi jetām}
  \item \textquotedblleft{}If I had money, I'd buy a boat\textquotedblright{} — \emph{ṭākā thākle āmi ekṭā noukā kintām}
  \item \textquotedblleft{}I used to go to the river every day\textquotedblright{} — \emph{āge āmi roj nôdīte jetām}
\end{itemize}

\begin{tcolorbox}[breakable,colback=teal!4,colframe=teal!35!black,title={Before you move on}]
\textquotedblleft{}I would go\textquotedblright{}? (\textbf{\bn{যেতাম}}.) \textquotedblleft{}He would go\textquotedblright{}? (\textbf{\bn{যেত}}.) What makes \textbf{\bn{যেতাম}} a habit? (A word like \textbf{\bn{রোজ}} or \textbf{\bn{প্রতিদিন}}, and a time in the past.)
\end{tcolorbox}
\section[\texorpdfstring{\bn{সঙ্গে} (sôngge)}{সঙ্গে (sôngge)}]{\bn{সঙ্গে} (sôngge) — with}

\label{lesson:BN-C275-songge}



\begin{quote}
Before the new one: say \textquotedblleft{}If I had time, I would go\textquotedblright{}. (\emph{jôdi sômôy thāktô, tāhôle āmi jetām}.) Then the line from your first meeting: \textbf{\bn{আলাপ} \bn{করে} \bn{ভালো} \bn{লাগলো}}.
\end{quote}

\subsection*{You'll want to know: \bn{সঙ্গে}}
\textbf{\bn{সঙ্গে}} — \emph{sôngge} — \textquotedblleft{}with\textquotedblright{}.

\begin{quote}
\textbf{\bn{আমি} \bn{বন্ধুর} \bn{সঙ্গে} \bn{সিনেমায়} \bn{গেলাম।}}
\end{quote}

\begin{quote}
\emph{āmi bôndhur sôngge sinemāy gelām.} — \textquotedblleft{}I went to the cinema with a friend.\textquotedblright{}
\end{quote}

The one you are with takes the \textbf{-\bn{র}} form: \textbf{\bn{বন্ধুর} \bn{সঙ্গে}}, \textbf{\bn{আমার} \bn{সঙ্গে}}, \textbf{\bn{তোমার} \bn{সঙ্গে}}. Now the greeting of chapter three can say who: \textbf{\bn{তোমার} \bn{সঙ্গে} \bn{আলাপ} \bn{করে} \bn{ভালো} \bn{লাগল}}, \textquotedblleft{}it was good to meet you\textquotedblright{}.

\begin{cousinweb}
\textbf{\bn{সঙ্গে}} (\emph{sôngge}) is usually read as \textbf{\bn{সঙ্গ}}, \textquotedblleft{}company\textquotedblright{}, with the locative \textbf{-\bn{এ}}: \textquotedblleft{}in the company of\textquotedblright{}. \textbf{\bn{সঙ্গ}} is a tatsama, Sanskrit \emph{saṅga}, which Monier-Williams gives as \textquotedblleft{}contact with\textquotedblright{} and \textquotedblleft{}association with\textquotedblright{}.
\end{cousinweb}

\subsection*{Your turn}
Say these aloud:

\begin{itemize}
\raggedright
  \item \emph{sôngge}
  \item \textquotedblleft{}I used to go to the river with Grandfather\textquotedblright{} — \emph{āge āmi dādur sôngge nôdīte jetām}
  \item \textquotedblleft{}I went to the cinema with a friend\textquotedblright{} — \emph{āmi bôndhur sôngge sinemāy gelām}
  \item \textquotedblleft{}It was good to meet you\textquotedblright{} — \emph{tomār sôngge ālāp kore bhālo lāglô}
\end{itemize}

\begin{tcolorbox}[breakable,colback=teal!4,colframe=teal!35!black,title={Before you move on}]
What does \textbf{\bn{সঙ্গে}} mean? (\textquotedblleft{}With\textquotedblright{}.) Which form does the person take? (The \textbf{-\bn{র}} form: \textbf{\bn{তোমার} \bn{সঙ্গে}}.)
\end{tcolorbox}
\section[\texorpdfstring{\bn{ঙ} (ng)}{ঙ (ng)}]{\bn{ঙ} (ng) — the letter inside \bn{সঙ্গে}}

\label{lesson:BN-W275-letter-uma}



\begin{quote}
Before the new one: say \textbf{\bn{সঙ্গে}}, \textquotedblleft{}with\textquotedblright{}. Then read \textbf{\bn{বাংলা}}: its \textbf{\bn{ং}} says \emph{ng} too.
\end{quote}

\subsection*{Script you'll notice: \bn{ঙ}}
\textbf{\bn{ঙ}} — \emph{ng}, as in \textquotedblleft{}sing\textquotedblright{}. It almost never starts a word. In \textbf{\bn{সঙ্গে}} it is joined to \textbf{\bn{গ}} in one conjunct, \textbf{\bn{ঙ্গ}}, \emph{ngg}, which does not look like the two letters set side by side: learn that shape as a whole.

In the print of this book, \textbf{\bn{ঙ}} has no headline. Its lower part is a large round bowl: its left arm ends high, at the upper left, in a slanted tip, and its right side ends in a short flat end at the right. At the top, in the middle, sits a small closed loop, with a short straight arm slanting up to the left from it. Below the loop, a short stem and a curve join it to the end of the bowl's right side.

\subsection*{Writing: \bn{ঙ} — copy the printed shape}
Put your pen beside \textbf{\bn{ঙ}} and copy the shape you can see, slowly and large. Check the copy against the print: the bowl and its tall left arm, the small loop and its arm, the short stem and curve.

\begin{quote}
This book does not yet tell you where to start \bn{ঙ} or which way the pen goes. That is not written down here until it can be written down with a source. Copying the printed shape needs no such source, and it is how the shape becomes yours in the meantime.
\end{quote}

\subsection*{Your turn}
\begin{itemize}
\raggedright
  \item \emph{Find:} the conjunct \bn{ঙ্গ} in \textbf{\bn{সঙ্গে}}
  \item \emph{Write it:} \bn{ঙ}, from the print, then \textbf{\bn{ঙ্গ}}
  \item \emph{Write it:} \textbf{\bn{সঙ্গে}}
  \item \emph{Write it:} \textbf{\bn{ঢোকা}}, from memory
\end{itemize}

\begin{tcolorbox}[breakable,colback=teal!4,colframe=teal!35!black,title={Before you move on}]
Which letter is inside the \textbf{\bn{ঙ্গ}} of \textbf{\bn{সঙ্গে}}? (\textbf{\bn{ঙ}}.) Does \textbf{\bn{ঙ}} have a headline? (No.) Which other sign says \emph{ng}? (\textbf{\bn{ং}}, in \textbf{\bn{বাংলা}}.)
\end{tcolorbox}
\section[(six lines)]{Six lines about a trip}

\label{lesson:BN-W275-six-lines}



\begin{quote}
The A2 message asked for a practical note to a named reader. This task asks for six lines, and they tell a story: a trip, by train, bus or boat, real or invented.

Pick one reader: Arun, \textbf{\bn{অরুণ}}, or Mira, \textbf{\bn{মীরা}}. Say \textbf{\bn{সঙ্গে}} and \textbf{\bn{নৌকা}} aloud; both may go in.
\end{quote}

\subsection*{What to know first}
A told story is not a list of events. It has a \textbf{scene}, in \textbf{-\bn{ছিল}-}, or with \textbf{\bn{যখন} … \bn{তখন}}. It has \textbf{events} in the plain past, kept in order with \textbf{\bn{প্রথমে}}, \textbf{\bn{তারপর}} and \textbf{\bn{শেষে}}, and joined with the \emph{-e} form. It has a \textbf{reason}, with \textbf{\bn{কারণ}}, \textbf{\bn{তাই}} or \textbf{-\bn{এর} \bn{জন্য}}, and a \textbf{feeling}, with \textbf{\bn{আমার} … \bn{লাগল}} or \textbf{\bn{ভয়} \bn{করছিল}}.

A good one ends by looking at what might have been: \textbf{\bn{যদি} …, \bn{তাহলে} … \bn{যেতাম}}.

\subsection*{Writing — connected composition}
On paper, in Bengali script, write \textbf{six lines} to your reader about the trip, one line for each of these:

\begin{enumerate}
\raggedright
  \item the scene: when it was, and what you were doing as it began, with \textbf{-\bn{ছিলাম}} or \textbf{\bn{যখন} … \bn{তখন}};
  \item the first thing that happened, with \textbf{\bn{প্রথমে}}, and who you went with, with \textbf{\bn{সঙ্গে}};
  \item what happened next, with \textbf{\bn{তারপর}} or an \emph{-e} form, and one thing that went wrong: \textbf{\bn{হারিয়ে} \bn{গেছে}} or \textbf{\bn{ছেড়ে} \bn{গেছে}}, or \textbf{\bn{হঠাৎ}};
  \item why, with \textbf{\bn{কারণ}}, \textbf{\bn{তাই}} or \textbf{-\bn{এর} \bn{জন্য}};
  \item how it ended and how you felt, with \textbf{\bn{শেষে}} and \textbf{\bn{আমার} … \bn{লাগল}} or \textbf{\bn{ভয়} \bn{করছিল}};
  \item what might have been, with \textbf{\bn{যদি} …, \bn{তাহলে} … -\bn{তাম}}, and how far you went, with \textbf{\bn{দূর}} or \textbf{\bn{থেকে} … \bn{পর্যন্ত}}.
\end{enumerate}

There is no model in this lesson. Use only words you already write. The work is choosing the right past for each line, and joining them.

\begin{tcolorbox}[breakable,colback=teal!4,colframe=teal!35!black,title={Before you move on}]
Now read your lines as you read Mira's Saturday. Mark every verb \textbf{s} for scene, \textbf{e} for event. Line one should be mostly \textbf{s}; lines two and three mostly \textbf{e}.

Then three checks. A past ending is \textbf{-\bn{ল}}, said \emph{-lô}: \textbf{\bn{গেল}}, \textbf{\bn{এল}}. \textbf{\bn{নৌকা}} has \textbf{\bn{ৌ}} around its \textbf{\bn{ন}}, and \textbf{\bn{সঙ্গে}} has \textbf{\bn{ঙ}} joined to \textbf{\bn{গ}}. \textbf{\bn{হঠাৎ}} ends in \textbf{\bn{ৎ}}.

Last, read line six alone. Do both halves of the if take \textbf{-\bn{ত}-}: \textbf{\bn{থাকত}}, \textbf{\bn{যেতাম}}?
\end{tcolorbox}
\section[(review)]{Review — it depends, and what would be}

\label{lesson:BN-R275-depends-and-would}



\begin{quote}
Say \textquotedblleft{}It depends on the weather\textquotedblright{}, then \textquotedblleft{}If I had time, I would go\textquotedblright{}. (\emph{eṭā ābôhāoyār opôr nirbhôr kôre}; \emph{jôdi sômôy thāktô, tāhôle āmi jetām}.)
\end{quote}

\subsection*{Your turn}
\begin{itemize}
\raggedright
  \item \emph{Say it:} Will we go tomorrow? It depends on the weather. If Mira doesn't come, we won't go.
  \item \emph{Say it:} It's up to you. If you come, I'll come with you.
  \item \emph{Say it:} If I had money, I would buy a boat.
  \item \emph{Say it:} In the old days I used to go to the river with Grandfather every day.
  \item \emph{Say it:} We went into the shop with Mira.
  \item \emph{Write it:} \textbf{\bn{সঙ্গে}}, then \textbf{\bn{ঢোকা}}
\end{itemize}

\begin{tcolorbox}[breakable,colback=teal!4,colframe=teal!35!black,title={Before you move on}]
\textquotedblleft{}It depends on\textquotedblright{}? (\textbf{-\bn{এর} \bn{ওপর} \bn{নির্ভর} \bn{করে}}.) \textquotedblleft{}I would go\textquotedblright{}? (\textbf{\bn{যেতাম}}.) \textquotedblleft{}I used to go\textquotedblright{}? (\textbf{\bn{রোজ} \bn{যেতাম}}.) Which letter is inside \textbf{\bn{সঙ্গে}}? (\textbf{\bn{ঙ}}.)
\end{tcolorbox}
